With the advent of the High-Luminosity Large Hadron Collider (HL-LHC), the ATLAS detector will face the unprecedented challenge of reconstructing particle trajectories in an environment featuring up to 200 simultaneous proton-proton interactions per bunch crossing. To meet these challenges, the ATLAS Collaboration has decided to adopt the ACTS (A Common Tracking Software) toolkit --- a community-driven project that provides experiment-independent tracking algorithms implemented in modern C++.

This thesis presents developments and performance evaluations of track reconstruction software optimized for the conditions after the High-Luminosity upgrade of the ATLAS experiment at CERN. The work detailed in this thesis contributes significantly to the ACTS codebase, enhancing its capabilities across a range of reconstruction tasks. These include improvements in track propagation, stepper mechanics, navigation logic, ambiguity resolution, and time-aware tracking. Special focus is given to developments that improve the computational efficiency and robustness of track finding under HL-LHC conditions.

In addition to software developments, this thesis includes comprehensive tracking performance studies with the OpenDataDetector, a prototype of an (HL-)LHC-like full-silicon inner tracking detector, as well as the ATLAS Inner Tracker. Results demonstrate the effectiveness of the new developments in high-pileup environments, demonstrating track finding efficiency, resolution, and computational performance.

This work establishes a robust foundation for the integration of ACTS into the ATLAS track reconstruction software and contributes to a growing toolkit that supports future high-energy physics experiments.
